\section{Introduction}\label{sec:intro}

\subsection{The problem}

Can a deterministic multitape Turing machine transpose a $d\times d$ matrix with $b(d)=\max(1,\clog d)$-bit entries, written
row by row, in time linear in the input length $L=d^2b(d)$? This is the statement~(T) of the first paper of this
series~\cite{FirstPaper2026}. By a reduction of Harvey and van der Hoeven~\cite[Cor.~6.9]{HvdH2025}, a linear-time
multiplication machine would give a linear-time transposer, so a proof of $\neg$(T) would give the first superlinear lower
bound for off-line integer multiplication on multitape Turing machines. The question is open, and this paper does not decide
it.

A machine is \emph{oblivious} if its head motion and the steps at which it appends output bits depend only on~$d$. Every machine can
be made oblivious at a logarithmic cost in time~\cite{PippengerFischer1979}, so oblivious machines are the natural first target. The
run of an oblivious machine is a fixed time--space network, and a linear-time oblivious transposer is a network code for the
transposition demands on that network. Adler, Harvey, Jain, Kleinberg and Lehman showed that if coding could not beat routing on
undirected networks, oblivious two-tape transposition would need time $\Omega(L\log d)$~\cite[\S6]{AHJKL2006}; Braverman and He, and
concurrently Zhang, Li and Li, have refuted that hypothesis~\cite{BravermanHe2026,ZhangLiLi2026}, so for oblivious machines the question
is genuinely one about codes. The series has reduced it to codes in two steps. The fourth paper reduced $\neg$(T) for oblivious machines
to depth reduction and a coding statement, the \emph{crux}~\cite[Theorem~D]{FourthPaper2026}. Lower bounds for several restricted
classes of labels suggested that depth reduction might fail~\cite{FourthPaper2026,EighthPaper2026,DirectedReciprocity2026,MergeDepth2026},
but the fourteenth paper proved it for every number of heads, so that the crux alone would exclude oblivious linear-time
transposers~\cite[Corollary~O]{BoundedStack2026}. The crux asks whether the \emph{cone relations} of a linear-time schedule, the
relations between each output bit and the input bits it can see without crossing a small revealed set, admit codes of rate close to~$1$.
Machine-free relations of this kind carry storage codes above rate one half~\cite{StorageCodes2026}, and binary codes of rate approaching
one exist on suitable separated relations~\cite{HighRate2026}; whether the relations of an actual linear-time schedule can carry them is
the question.

This paper attacks the crux from the side of the schedule. Instead of a coding bound for all relations, we ask which schedules make the
crux \emph{provable}. We show that it holds, with an explicit deficit, for every schedule in which a small enough set of steps separates a
large set of outputs from the \emph{macroscopic crossings} of the tapes. An oblivious linear-time transposer must therefore route a linear
amount of information through macroscopic crossings, and we describe where it must do so.

\subsection{Results}

The run of a machine defines a \emph{dependency DAG} on its steps $0,\dots,T-1$: every step depends on the step before it, which
carries the state, and on the previous visit of every work head to the cell it is scanning, which carries that cell's
contents. A \emph{revisit} of a work tape is a pair $u<t$ of consecutive visits of its head to one cell between which the head leaves the
cell; its \emph{excursion} is the farthest the head travels from the cell before returning. Two revisits $u<t$ and $u'<t'$ of the same tape
\emph{cross} if $u<u'<t<t'$. A unit-speed head makes crossing revisits only when it reuses a stretch of tape like a queue, across two
turns: write a stretch moving right, return, and read it again moving right. Reuse like a stack, across one turn, makes only nested
revisits. Given $\theta>0$, a step is a \emph{macroscopic crossing step} if it is the earlier end of a revisit of excursion more than
$\theta L$ that crosses another revisit of excursion more than $\theta L$ on the same tape. A trace of a machine is \emph{$\eta$-heavy} if at
least a $2^{-\eta L}$ fraction of the inputs follow it; an oblivious machine has one trace, and it is $0$-heavy. \Cref{sec:prelim} states
these definitions exactly.

\begin{itemize}[leftmargin=1.5em]
\item \textbf{Theorem~X, macroscopic crossings} (\cref{thm:X}). Let $M$ be an oblivious machine with $k$ work tapes that
transposes $d\times d$ matrices within $C(L+1)$ steps, and let $\alpha>0$. There are explicit $\theta>0$, $c>0$ and $d_0$,
depending only on $k$, the alphabet and state counts of~$M$, $C$ and~$\alpha$, such that for every $d\ge d_0$ the following
holds. For every set $V_0$ of at least $\alpha L$ output positions, every set of steps that meets every dependency path from
the output steps of $V_0$ to the macroscopic crossing steps has at least $c\abs{V_0}$ elements. By Menger's theorem, there are then
$c\abs{V_0}$ vertex-disjoint dependency paths from these outputs into macroscopic crossings.
\item \textbf{Theorem~D, a dichotomy for all traces} (\cref{thm:D}). Theorem~X is the transposer case of a statement about
arbitrary linear-time traces, with no hypothesis on what the machine computes: for every output set of size at least $\alpha L$,
either the macroscopic crossings are reached with linear expansion, or a cheap separator leaves a linear-size acyclic set in the cone
relation. A separator cheaper than the acyclic set it leaves, at a radius at least the genie depth, is impossible for a transposer
(\cref{lem:acyclic-cost}).
\item \textbf{Theorem~S, skeletons} (\cref{thm:S}). At scale $\theta L/4$ the walk of each work tape of a linear-time run is a
word of at most $\Mmax=\lfloor16C/\theta\rfloor$ steps $\pm1$, so the coarse shapes of all runs come from an explicit finite
family of $(2^{\Mmax+1}-1)^k$ tuples, independent of~$d$. For an oblivious linear-time transposer the shape is
\emph{crossing-capable} (some tape turns at least twice at this scale with overlapping legs), and for every large output set the
expansion of Theorem~X is carried by the crossing steps of a \emph{single coarsely monotone leg of a single tape}. Along infinitely many
sizes, one fixed shape, tape and leg serve for the set of all input positions.
\item \textbf{Theorem~H, two long-range heads} (\cref{thm:H}). No machine that transposes in linear time with an $\eta$-heavy trace, for a
small fixed $\eta$, has on that trace one unrestricted work head while every other work head keeps its revisits within a small fixed
fraction of~$L$, whatever the input head does. Every such transposer, in particular every oblivious one, therefore has at least two work
heads that return to cells after moving a fixed fraction of $L$ away; and for each work head, some set containing a fixed fraction of its
outputs reaches long revisits off that head with linear expansion.
\item \textbf{Theorem~E, adaptive machines} (\cref{thm:E}). The conclusion of Theorem~X holds for every $\eta$-heavy trace of an arbitrary
linear-time transposer, for an explicit $\eta>0$. It holds on all heavy traces of the bit simulator of the fifteenth
paper~\cite[Theorem~S]{HighRate2026} applied to any linear-time transposer; that simulator has the same tapes and work alphabet, runs in
linear time, and its heavy traces carry all but a $2^{-\eta L}$ fraction of the inputs. Transferring the conclusion back to the original
machine's own dependency DAG is open; every statement in this paper about the dependency DAG of a general machine assumes a heavy trace or
names the simulator.
\item \textbf{Lemma~R, rows to columns} (\cref{lem:R}). For every machine that transposes at size $d$, every $\eta$-heavy trace, and every two
sets $A_0,A_1$ of rows, every set of steps that separates the input reads of the rows in $A_0$ from the output steps of the entries that go
from rows $A_0$ to output rows $A_1$ has at least $(\abs{A_0}\abs{A_1}b-\eta L)/(\wval+1)$ elements, where $\wval$ is the number of bits of
one step's value.
\item \textbf{Class exclusions} (\cref{sec:classes}). Machines whose heads make boundedly many macroscopic phases, oblivious or adaptive with
sublinear jitter, are excluded by an explicit native deficit (Theorems~B and~J); a fixed native deficit suffices in place of the half-rate
hypothesis of the fourteenth paper's low-adaptivity corollary (Theorem~G); and machines with no free work head, bounded nesting load and
boundedly many scales are excluded by a scale-gap pigeonhole (Theorem~Q). \Cref{app:catalogue} lists explicit schedule families excluded
by these theorems and by certificates, with proof sketches.
\end{itemize}

\paragraph{What remains.} Explicit oblivious linear-time schedules, realized in our experiments by machines that compute other functions,
meet the necessary conditions above in every test we could run (numerically, at $d\le128$; \cref{sec:construction}). So these conditions
are not known to decide the question. Counting cuts alone cannot go further: one set of steps, the storage live across the last input read
together with the earlier output steps, separates the inputs from the outputs for every pair of row sets at once (\cref{prop:onecut}).
Outside restricted schedule classes (\cref{cor:scale-separated}), the only coding-proof logarithmic factor known on tapes comes from counting
passes (Theorem~P of the first paper), and \cref{sec:open} states the multi-scale question it leaves open.

\paragraph{Constants.} The constants of Theorems~X, D and~S are explicit, positive and astronomically small. For two work tapes, time
constant $C=1$ and $\alpha=1$, the expansion constant of Theorem~X is $\rho_0/(4\wval)$ with $\rho_0\approx10^{-105}$, and the threshold $\theta$
is smaller still; \cref{sec:crossings} gives the formulas. The constants of Theorem~H are polynomial in the time constant and the value width,
and those of Theorem~Q come from a compactness argument and are not explicit. The small constants come from the proofs: on the schedules we
measured, the expansion ratio was exactly~$1$ at every tested threshold up to $L/2$ (numerically, at $d\le128$; \cref{rem:numerics-crossings}).

\subsection{How to read the paper and its formal proofs}

Every numbered result states its status. \textsf{Lean:} lines name the declarations of the accompanying Lean~4 library that
prove it, in the file named at the right; the Lean statement is authoritative and the text gives a faithful outline of the formal
proof. A result marked \textsf{Not formalized} has a complete written proof in the text. \textsf{Cited:} lines name results of
earlier papers of the series used without proof, and say whether they are Lean-verified there. Numerical observations appear
only in remarks or parentheses labelled as numerical and are never used in proofs. The formal library depends only on the axioms
\lean{propext}, \lean{Classical.choice} and \lean{Quot.sound} (\cref{app:verification}).

The formalization is in the frozen machine model of the first paper~\cite{FirstPaper2026}, through the traces, dependency DAGs
and index codes of the eleventh paper~\cite{MergeDepth2026}, the genies of the fourteenth~\cite{BoundedStack2026} and the
simulator of the fifteenth~\cite{HighRate2026}. Not formalized: \cref{prop:onecut} and \cref{cor:scale-separated}, which
have complete proofs in the text; the remarks marked as proof sketches; and the exclusions of the individual schedule families of
\cref{app:catalogue}, which are not numbered results and are given as proof sketches over the formal class theorems and certificate
lemmas.

\subsection{Related work}

\emph{Lower bounds on tapes.} Superlinear lower bounds for explicit problems are known for one work tape: Hennie's crossing
sequences~\cite{Hennie1965}, Maass's combinatorial arguments~\cite{Maass1985}, and for transposition itself the bounds of
Dietzfelbinger, Maass and Schnitger~\cite{DMS1991} and Dietzfelbinger and Maass~\cite{DietzfelbingerMaass1993} in the classical one-tape
models. With two or more work tapes the only superlinear lower bounds known are those that follow from the separation of deterministic
and nondeterministic linear time~\cite{PPST1983}, for problems that are hard for nondeterministic linear time~\cite{Grandjean1988}; none is
known for off-line multiplication, sorting, the discrete Fourier transform or transposition~\cite[\S1.4]{HvdH2025}. The simulation of time
by space~\cite{HPV1977} also rests on the structure of multitape computation graphs. For transposition, moving the entries as indivisible
items along a fixed pattern needs time $\Omega(L\log d)$~\cite{Stoss1973}, and similar bounds hold for indivisible records in external
memory~\cite{AggarwalVitter1988}; Kirchherr proved an $\Omega(\log d)$ bound on reversals for conservative machines~\cite{Kirchherr1988}, and
the first paper of this series an $\Omega(\log d)$ bound on monotone passes for arbitrary computation, with coding
allowed~\cite[Theorems~P and~P$'$]{FirstPaper2026}, extended to adaptive machines without a runtime assumption
in~\cite{PassesMemo2026}. The stream lower bounds of Grohe, Hernich and Schweikardt~\cite{GroheSchweikardt2005,GHS2009} bound reversals
for sorting-type problems with small internal memory. A linear-time machine, however, may reverse its heads $\Theta(L)$ times at small
scales, so pass and reversal bounds apply to it only after its small-scale motion is absorbed into the finite control.

\emph{Macroscopic motion.} The first paper's Theorem~M*~\cite[Theorems~7.8 and~7.9]{FirstPaper2026} passes to the limit of rescaled head
trajectories and shows that the heads of any linear-time transposer, oblivious or adaptive, satisfy an integral inequality forcing
returns to macroscopically earlier positions. The fourth paper deduced, by a written proof, that in every macroscopic limit of a linear-time
transposer at least two work-head limits are non-monotone~\cite[Corollaries~3.18 and~3.24]{FourthPaper2026}. Theorem~H(b) is a
quantitative form of this at a single size, with explicit constants and on heavy traces. Theorem~X asks for more than non-monotone motion: a
run whose heads reuse their tapes only like stacks has no macroscopic crossing steps, so Theorem~X excludes it whatever its macroscopic limit.

\emph{Network and index coding.} The time--space view of oblivious machines and the conditional lower bound are due to Adler et
al.~\cite{AHJKL2006}; the multiple-unicast conjecture was refuted by Braverman and He~\cite{BravermanHe2026} and, concurrently, by Zhang, Li
and Li~\cite{ZhangLiLi2026}. Our exclusions use the acyclic-set lower bound of index coding~\cite{BBJK2011}, in the machine form of the
eleventh paper~\cite{MergeDepth2026}.

\emph{Upper bounds.} Recursive block transposition takes time $O(L\log d)$. An unrefereed preprint claims time
$O(L(\log d)^{1-\kappa})$ for a small $\kappa>0$ through an address-interchange primitive that mixes data by exclusive
or~\cite{OpenAI2026mult}; if it is correct, coding already beats routing on tapes by a factor $(\log d)^{\kappa}$, and any
proof of $\neg$(T) must allow such gains.

\subsection{Organization}

\Cref{sec:prelim} sets out the model, the dependency DAG, revisits and crossings, cuts and cone relations. \Cref{sec:certificates}
gives the certificate method: an acyclic set of linear size in a cone relation, possibly after a cheap cut, excludes linear-time
transposition. \Cref{sec:heads} proves Theorem~H, \cref{sec:crossings} Theorems~D and~X, and \cref{sec:skeletons} Theorem~S.
\Cref{sec:demand} proves Lemma~R and the colouring criterion built on it. \Cref{sec:adaptive} treats adaptive machines,
\cref{sec:classes} the class exclusions. \Cref{sec:construction} turns the results around into a specification for a linear-time
transposer and describes the schedules that meet it, and \cref{sec:open} lists the open problems that would decide the question.
\Cref{app:catalogue} catalogues the excluded schedule families, \cref{app:leanindex} indexes the Lean declarations, and
\cref{app:verification} explains how to build and check the formalization.
